% =============================================================================
\chapter{Réalisation — phases D à H}\label{chap:phases-dh}
% =============================================================================

% -----------------------------------------------------------------------------
\section{Phase D — Automatisation avec Ansible}\label{sec:phase-d}

\fiche{Fournir des playbooks Linux et Windows reproductibles et idempotents.}
      {Kenzi Bendjelloul}{19/10/2026}{\fait}

\subsection{Réalisation}

Ansible est installé sur OPS01 avec les collections du dépôt (\code{ansible/requirements.yml}).
Un compte de service \code{ansible} dédié (clé ed25519, sudo sans mot de passe validé par
\code{visudo -c}) administre LINUX01 par SSH ; WIN-DC01 est administré par WinRM (HTTP 5985, NTLM)
avec le mot de passe Administrateur lu dans un \textbf{vault chiffré AES256}.

\begin{table}[H]
\centering
\caption{Playbooks du dépôt}\label{tab:playbooks}
\begin{tabularx}{\linewidth}{@{}P{3.4cm}P{1.9cm}LP{3.6cm}@{}}
\toprule
\textbf{Playbook} & \textbf{Cible} & \textbf{Actions} & \textbf{Vérifications intégrées} \\
\midrule
\code{windows_adds.yml} & WIN-DC01 & Rôle AD DS, forêt, NTP, redirecteurs, client DNS & Services AD, état du domaine \\
\code{windows_base.yml} & WIN-DC01 & OU, groupes, comptes, partage SMB et ACL NTFS, zone inverse, enregistrements A/PTR, windows\_exporter 0.31.8, pare-feu 9182 limité à OPS01 & Services critiques démarrés, \code{/metrics} joignable \\
\code{linux_base.yml} & LINUX01 & Paquets, comptes, droits, site nginx (modèle + rechargement), \code{/data} par UUID, pare-feu UFW, node\_exporter, contrôle planifié (séance 6) & HTTP 200, contenu publié, \code{/metrics} \\
\code{ops01_monitoring.yml} & OPS01 & Prometheus, règles, Grafana provisionné & Toutes les cibles \code{up}, compte admin Grafana \\
\code{linux_maintenance.yml} & Linux & Maintenance (phase E) & Pré- et post-contrôles \\
\code{windows_maintenance.yml} & WIN-DC01 & Windows Update (phase E) & Pré- et post-contrôles \\
\code{controle_sante.yml} & 3 serveurs & Aucune : contrôle de santé en lecture seule (séance 6) & Tableau OK / ATTENTION / ECHEC \\
\code{linux_quotas.yml} & LINUX01 & Quotas sur le disque de test (séance 6) & Garde-fou sur le numéro de série \\
\bottomrule
\end{tabularx}
\end{table}

Deux contrôles ont été ajoutés le 20/10 par le chef de projet : chaque playbook Windows et la
supervision commencent par vérifier que les variables du vault sont disponibles, \emph{avant toute
modification} ; le script \code{scripts/verif-idempotence.sh} enchaîne deux exécutions et renvoie
0 si le second passage est à \code{changed=0}, 1 s'il reste des changements, 2 en cas d'échec.

\subsection{Résultats}

\begin{table}[H]
\centering
\caption{Phase D — vérifications et résultats}\label{tab:phase-d}
\begin{tabularx}{\linewidth}{@{}P{3cm}LP{2.6cm}@{}}
\toprule
\textbf{Vérification} & \textbf{Résultat observé} & \textbf{Preuve} \\
\midrule
Inventaire et connexions & 3 hôtes, 3 groupes ; \code{ping: pong} et \code{win_ping: pong} & \preuve{P-30} \\
1\textsuperscript{er} passage & Linux : 22 tâches OK dont 10 changements ; Windows : 21 tâches OK dont 12 changements ; aucun échec & \preuve{P-31}, \preuve{P-33} \\
2\textsuperscript{e} passage (idempotence) & Linux \code{changed=0}, Windows \code{changed=0} & \preuve{P-32}, \preuve{P-34} \\
Script d'idempotence & Code 0 sur les playbooks réels, 1 sur un playbook non idempotent, 2 sur un playbook en échec & \preuve{P-61}, \preuve{P-62} \\
Contrôle du vault & Arrêt avant toute tâche si le vault est absent ou invalide & \preuve{P-62} \\
\bottomrule
\end{tabularx}
\end{table}

\subsection{Difficultés et décisions}

\begin{table}[H]
\centering
\caption{Phase D — difficultés rencontrées}\label{tab:phase-d-diff}
\begin{tabularx}{\linewidth}{@{}P{3.8cm}P{4cm}L@{}}
\toprule
\textbf{Problème} & \textbf{Cause} & \textbf{Correction et résultat} \\
\midrule
L'option MSI de windows\_exporter ouvre le port 9182 à tous & Comportement de l'installeur & Option retirée, règle dédiée limitée à OPS01 \\
Le PTR du DC repasse en \code{changed} & Le DC réenregistre lui-même son PTR (dynamique, TTL 20 min) ; le playbook le réécrivait en statique : chacun écrasait l'autre & Le playbook déclenche seulement l'enregistrement du DC s'il manque : \code{changed=0} même après réenregistrement forcé (\preuve{P-36b}, \preuve{P-37}) \\
Test négatif du vault exécuté réellement & \code{--vault-password-file} s'ajoute au fichier configuré au lieu de le remplacer & Impact vérifié (valeur inchangée) ; test refait avec \code{ANSIBLE_VAULT_PASSWORD_FILE} ; défaut d'idempotence révélé puis corrigé (\preuve{P-62}, \preuve{P-15}) \\
1\textsuperscript{er} pipeline CI en échec & Dashboard absent, ligne trop longue, 7 remarques ansible-lint & Corrigé ; mêmes outils installés en local pour valider avant de pousser \\
\bottomrule
\end{tabularx}
\end{table}

% -----------------------------------------------------------------------------
\section{Phase E — Maintenance et mises à jour}\label{sec:phase-e}

\fiche{Formaliser et tester une maintenance sûre et traçable sur les deux systèmes.}
      {MIR Mahan, Jolan Allen}{19/10/2026}{\fait}

La procédure (\code{docs/procedure-maintenance.md}) fixe la fenêtre, l'ordre (LINUX01, puis
OPS01, puis WIN-DC01 en dernier car tout le laboratoire en dépend), les précautions et le retour
arrière. Elle est outillée par deux playbooks qui \textbf{refusent de continuer} si un pré-contrôle
échoue.

\begin{table}[H]
\centering
\caption{Étapes de la maintenance outillée}\label{tab:maintenance}
\begin{tabularx}{\linewidth}{@{}P{2.4cm}LL@{}}
\toprule
\textbf{Étape} & \textbf{Linux (\code{linux_maintenance.yml})} & \textbf{Windows (\code{windows_maintenance.yml})} \\
\midrule
Retour arrière & Snapshot \code{s04_avant_maintenance_e} & Snapshot \code{s04_avant_maj_windows} \\
Pré-contrôles & Systèmes de fichiers < 85 \% ; services attendus actifs & Volumes fixes < 85 \% ; services AD démarrés ; aucun redémarrage en attente \\
Mises à jour & \code{unattended-upgrades} (sécurité quotidienne) ; \code{apt upgrade} sans suppression & Recherche, puis installation seulement avec \code{-e maintenance_apply=true} \\
Redémarrage & Seulement s'il est requis \textbf{et} autorisé (\code{-e allow_reboot=true}) & Seulement avec \code{-e allow_reboot=true} \\
Post-contrôles & Services actifs ; site HTTP 200 & Services critiques ; domaine et DNS ; historique Windows Update \\
Trace & Rapport par serveur dans \code{~/infraready-logs/} & Idem + journal Windows Update \\
\bottomrule
\end{tabularx}
\end{table}

\begin{table}[H]
\centering
\caption{Phase E — vérifications et résultats}\label{tab:phase-e}
\begin{tabularx}{\linewidth}{@{}P{2.6cm}LP{2.7cm}@{}}
\toprule
\textbf{Vérification} & \textbf{Résultat observé} & \textbf{Preuve} \\
\midrule
Mise à jour Linux & OPS01 mis à jour ; aucun redémarrage requis ; mises à jour de sécurité automatiques actives & \preuve{P-40}, \preuve{P-43} \\
Mise à jour Windows & 5 mises à jour installées en 1 h 13 avec redémarrage du DC (cumulative KB5126646, .NET, MSRT, Defender) ; passage complet après exclusion de KB5007651 & \preuve{P-41}, \preuve{P-42}, \preuve{P-42c} \\
Capacité disque & Tous les volumes < 85 \% ; \code{C:} de 24,2 \% à 37,2 \% après la cumulative & \preuve{P-40b}, \preuve{P-42c} \\
Services & Tous actifs avant et après ; site HTTP 200 & \preuve{P-40}, \preuve{P-40b} \\
\bottomrule
\end{tabularx}
\end{table}

\attention{Difficulté : boucle Windows Update sur KB5007651}{%
L'installation se termine en échec après 1 h 13 alors que cinq mises à jour sont installées : la
plateforme Windows Security KB5007651 est marquée « installée » trois fois de suite mais toujours
reproposée, et le module \code{win_updates} arrête la tâche. Diagnostic par l'historique Windows
Update, le journal Système et \code{Get-MpComputerStatus} (plateforme bien à jour). Décision :
exclusion de KB5007651 (\code{reject_list}), à suivre manuellement ; passage suivant complet,
\code{failed=0} (\preuve{P-42} → \preuve{P-42b} → \preuve{P-42c}). Autre difficulté corrigée : un
lecteur CD vide était vu comme un volume plein à 100 \% (\preuve{P-41a}).}

% -----------------------------------------------------------------------------
\section{Phase F — Supervision avec Prometheus et Grafana}\label{sec:phase-f}

\fiche{Collecter et visualiser les métriques réelles des serveurs et démontrer une alerte.}
      {MIR Mahan, Jolan Allen}{19/10/2026}{\fait}

\subsection{Réalisation}

Tout est déployé par \code{ops01_monitoring.yml} à partir des fichiers versionnés dans
\code{monitoring/} : configuration et règles de Prometheus (validées par \code{promtool} avant tout
redémarrage), source de données et tableau de bord Grafana provisionnés, mot de passe
administrateur Grafana lu dans le vault, accès anonyme en lecture seule. Les exporters sont
installés par \code{linux_base.yml} et \code{windows_base.yml} ; leurs ports ne sont ouverts qu'à
OPS01 (figure~\ref{fig:flux-f}).

\illustration{00-flux-metriques-phase-f.pdf}{0.95}{Origine, ports et cheminement des métriques (phase F)}{fig:flux-f}

\begin{table}[H]
\centering
\caption{Règles d'alerte (5 en phase F, 7 après la séance 6)}\label{tab:regles}
\begin{tabularx}{\linewidth}{@{}P{4.3cm}LP{1.2cm}P{1.6cm}@{}}
\toprule
\textbf{Alerte} & \textbf{Condition} & \textbf{Durée} & \textbf{Sévérité} \\
\midrule
\code{InstanceDown} & Cible Linux ou Windows injoignable (\code{up == 0}) & 1 min & critical \\
\code{ServiceWindowsCritiqueArrete} & Service DNS, NTDS, Netlogon ou Kdc arrêté & 1 min & critical \\
\code{CpuEleve} & CPU > 80 \% & 2 min & warning \\
\code{MemoireElevee} & Mémoire > 90 \% & 5 min & warning \\
\code{DisquePresquePlein} & Volume > 85 \% & 5 min & warning \\
\code{SiteWebIndisponible} (séance 6) & Contrôle planifié du site en échec & 2 min & critical \\
\code{ControlePlanifieEnRetard} (séance 6) & Contrôle planifié non exécuté depuis 5 min & 5 min & warning \\
\bottomrule
\end{tabularx}
\end{table}

\subsection{Résultats}

\illustration{P-51-prometheus-targets.png}{0.9}{Cibles Prometheus : 4/4 \code{up} (\preuve{P-51})}{fig:p51}

\begin{table}[H]
\centering
\caption{Phase F — vérifications et résultats}\label{tab:phase-f}
\begin{tabularx}{\linewidth}{@{}P{2.6cm}LP{2.6cm}@{}}
\toprule
\textbf{Vérification} & \textbf{Résultat observé} & \textbf{Preuve} \\
\midrule
Collecte Linux & LINUX01 \code{up} (1 709 métriques) ; CPU 0,2 \%, mémoire 11,3 \%, \code{/} 29,1 \% & \preuve{P-51}, \preuve{P-52} \\
Collecte Windows & WIN-DC01 \code{up} (3 263 métriques) ; CPU 1,0 \%, mémoire 26,7 \%, \code{C:} 24,2 \% ; services AD à 1 & \preuve{P-51}, \preuve{P-52} \\
Tableau de bord & 3 serveurs « EN LIGNE », CPU, mémoire, stockage, alertes ; version 2 avec filtre d'hôte (PR \#3) & \preuve{P-53}, \preuve{P-56} \\
Alerte démontrée & \code{InstanceDown} pending à +10 s, firing à +60 s, résolue 16 s après le retour du service & \preuve{P-54}, \preuve{P-55} \\
Idempotence & \code{ops01_monitoring.yml} : \code{changed=0} au 2\textsuperscript{e} passage & \preuve{P-50b} \\
\bottomrule
\end{tabularx}
\end{table}

\illustration{P-55-grafana-dashboard-alerte-linux01-hors-ligne.png}{0.95}{Dashboard Grafana (version 1) pendant la démonstration : LINUX01 hors ligne et alerte active (\preuve{P-55})}{fig:p55}

\attention{Difficulté : mot de passe Grafana tronqué}{%
Après le premier passage, le compte administrateur Grafana était refusé : le mot de passe contenait
\og \# \fg{}, lu comme un début de commentaire dans \code{grafana.ini}. Correction : mot de passe sans
caractère ambigu, contrôle \code{assert} dans le playbook et vérification sur l'API
d'administration (\preuve{P-50a} → \preuve{P-50}).}

% -----------------------------------------------------------------------------
\section{Phase G — GitHub et intégration continue}\label{sec:phase-g}

\fiche{Versionner les livrables et valider le dépôt automatiquement.}
      {MIR Mahan, Jolan Allen}{19--23/10/2026}{\fait}

Le dépôt \code{jolanallen/InfraReady} regroupe playbooks, scripts, supervision, modèles
d'installation, documentation et preuves (annexe~\ref{annexe:arborescence}). Le workflow
\code{.github/workflows/validation.yml} s'exécute à chaque push et à chaque demande de fusion.

\begin{table}[H]
\centering
\caption{Jobs du pipeline GitHub Actions}\label{tab:ci}
\begin{tabularx}{\linewidth}{@{}P{3.6cm}L@{}}
\toprule
\textbf{Job} & \textbf{Contrôles} \\
\midrule
Arborescence attendue & Présence des livrables (README, inventaire, playbooks, supervision) \\
YAML + Ansible & \code{yamllint --strict}, \code{ansible-playbook --syntax-check} de chaque playbook, \code{ansible-lint} \\
Prometheus + Grafana & \code{promtool check config} (configuration et règles), JSON du dashboard valide \\
Aucun secret en clair & gitleaks sur tout l'historique ; fichiers vault obligatoirement chiffrés \\
\bottomrule
\end{tabularx}
\end{table}

\begin{itemize}
  \item Premier run en échec (dashboard absent, lint) puis vert sur le dépôt de travail
        (\preuve{P-60a}) ; historique migré le 20/10 vers le dépôt de l'équipe, run vert
        (\preuve{P-60}).
  \item Depuis le 20/10, toute modification passe par une branche et une demande de fusion
        (16 au 23/10), relue et fusionnée par le chef de projet ; la CI est verte sur chacune
        (\preuve{P-61} à \preuve{P-64}).
  \item \textbf{Limite assumée} : la CI valide les fichiers mais ne déploie rien ; un runner hébergé par
        GitHub ne peut pas joindre le réseau privé du laboratoire. Le déploiement se fait depuis OPS01.
\end{itemize}

% -----------------------------------------------------------------------------
\section{Phase H — Diagnostic d'incident}\label{sec:phase-h}

\fiche{Reproduire un incident contrôlé, en trouver la cause et prouver le retour à la normale.}
      {MIR Mahan, Jolan Allen}{19/10/2026}{\fait}

\paragraph{Incident INC-2026-10-19-01.} À 19:27, la directive \code{root} du site nginx de
LINUX01 est modifiée vers un répertoire inexistant (\og mise en production \fg{} simulée,
sauvegarde préalable). \code{nginx -t} ne vérifie pas l'existence du répertoire : le rechargement
réussit et le site répond 404 pour toutes les pages.

\illustration{00-incident-phase-h.pdf}{0.42}{Chronologie de l'incident de la phase H}{fig:incident-h}

\begin{table}[H]
\centering
\caption{Phase H — démarche et résultats}\label{tab:phase-h}
\begin{tabularx}{\linewidth}{@{}P{2.4cm}LP{2.2cm}@{}}
\toprule
\textbf{Étape} & \textbf{Constat} & \textbf{Preuve} \\
\midrule
Symptôme & HTTP 404 depuis OPS01 par adresse et par nom ; VM joignable, port 80 ouvert & \preuve{P-70} \\
Diagnostic & Trois hypothèses écartées (service, pare-feu, contenu) ; \code{nginx -T} révèle le mauvais \code{root} ; \code{stat}, \code{diff} et le journal sudo datent et attribuent le changement & \preuve{P-71} \\
Correction & Rejeu de la configuration de référence par \code{linux_base.yml} plutôt qu'une correction manuelle & \preuve{P-72a}, \preuve{P-72} \\
Retour à la normale & HTTP 200 par adresse et par nom ; 2\textsuperscript{e} passage \code{changed=0} & \preuve{P-73} \\
\bottomrule
\end{tabularx}
\end{table}

\alerte{Défaut révélé dans le playbook}{%
La première correction restaurait bien le fichier mais le site restait en 404 : la vérification
HTTP échouait avant la fin du play, et Ansible n'exécute les gestionnaires (ici, le rechargement de
nginx) qu'en fin de play. Correction : \code{meta: flush_handlers} avant les vérifications.
L'incident a été reproduit une seconde fois pour prouver que le playbook corrigé le résout en un
seul passage. Rapport complet : \code{docs/rapport-incident.md}.}
